\section{噪声分类法：把不确定性拆成可建模组件}

\subsection{三层噪声结构}
Sida 把噪声源拆为可以落地操作的几层：task/data 层、inference/protocol 层、evaluation/judge 层。这种分层的价值在于，它把 ``噪声'' 从模糊概念变成可诊断对象。团队在调研 benchmark 波动时，常见失败是把所有误差都归因到模型；分层框架要求先定位误差来自哪一层。

\begin{table}[H]
\centering
\begin{tabular}{p{3.2cm}p{4.2cm}p{4.2cm}}
\toprule
\textbf{噪声层} & \textbf{典型来源} & \textbf{可控手段} \\
\midrule
Task/Data & 数据污染、难度漂移、样本选择偏差 & 数据去污染、难度分桶、分布追踪 \\
Inference/Protocol & prompt 模板、解码策略、上下文拼接方式 & 协议固定、seed 控制、A/B 模板对照 \\
Evaluation/Judge & rubric 模糊、judge 偏置、评审一致性不足 & 双评审、仲裁机制、置信区间汇报 \\
\bottomrule
\end{tabular}
\caption{可预测噪声的分层拆解框架}
\end{table}

\subsection{``可预测'' 的含义}
可预测不等于可消除。实际系统里，很多噪声不会消失，但可以被估计、被约束、被报告。比如当你知道某个 benchmark 在固定配置下的 run-to-run standard deviation 是 $0.8$，那么 $0.3$ 的改进就不应被当作 ``胜出''。这个思想本质是 measurement science，而不是 leaderboard engineering。

\begin{importantbox}{错误的报告方式}
只报单点分数，不报 variance；只报最佳 run，不报 run distribution；只报总体平均，不报分桶结果。
\end{importantbox}

\subsection{噪声与 ``可比性''}
社区经常做 cross-paper 对比，但如果协议不同、模板不同、judge 不同，分数本身不可比。Sida 的强调点是：先建立 comparability，再谈 superiority。否则论文 A 比论文 B 高 1 分，可能只是评测协议不同，而不是模型变强。

\begin{warningbox}{排行榜比较的常见误区}
当两个系统不共享同一评测协议时，排序信息的可信度会显著下降。盲目比较只会制造伪进展。
\end{warningbox}

\subsection{本章小结}
可预测噪声的关键不是 ``把噪声归零''，而是建立统一分层框架，让每个分数都能被解释。分层越清晰，后续统计估计和协议治理越容易落地。

